\documentclass[11pt]{article}

\usepackage[margin=1.15in]{geometry}
\usepackage[T1]{fontenc}
\usepackage{lmodern}
\usepackage{amsmath,amssymb}
\usepackage{booktabs}
\usepackage{microtype}
\usepackage{enumitem}
\usepackage{fancyvrb}
\usepackage[colorlinks=true,linkcolor=blue!45!black,urlcolor=blue!45!black,
            citecolor=blue!45!black]{hyperref}
\usepackage{titlesec}

\titleformat{\section}{\large\bfseries}{\thesection}{0.7em}{}
\titleformat{\subsection}{\normalsize\bfseries}{\thesubsection}{0.7em}{}

\setlength{\parskip}{0.35em}
\setlength{\parindent}{0pt}

\newcommand{\stcmd}[1]{\texttt{#1}}
\DefineVerbatimEnvironment{stata}{Verbatim}
  {fontsize=\small,frame=leftline,framerule=0.8pt,rulecolor=\color{black!25},
   xleftmargin=6pt,framesep=8pt}
\usepackage{xcolor}

\title{\textbf{copulaendog}: instrument-free copula corrections\\
       for endogenous regressors in Stata and Python}
\author{Ported from the R implementations of\\
        Rouven E.\ Haschka and Ashwin Malshe}
\date{}

\begin{document}
\maketitle

\begin{abstract}
\noindent
\stcmd{copulaendog} corrects endogenous regressors in a linear model when no
instrument is available. The dependence between the endogenous regressor and
the structural error is modelled with a Gaussian copula, and the resulting
control function is added to the regression. Five cross-sectional estimators
are implemented behind one interface: Park and Gupta (2012), the two-stage
copula generated regressor approach of Yang, Qian and Xie (2025), the
between-regressor correlation variant of Haschka (2025a), the
nonlinear-transformation approach of Breitung, Mayer and Wied (2024), and the
adjusted estimator of Liengaard et al.\ (2025). Seven marginal CDF estimators
are available for the copula transformation, inference is by pairs bootstrap,
and a \stcmd{validity} report walks the identification requirements the
literature states. This document describes the Stata command
\stcmd{copulaendog} and the companion Python package \stcmd{copulaendog}, both
ported from the R reference implementation of Haschka and from the
\stcmd{endogCopula} package of Malshe.
\end{abstract}

\section{Credit}

This software is a port. The econometrics, the design of the interface, and
the reference code are the work of two authors, and neither the Stata command
nor the Python package would exist without them:

\begin{itemize}[leftmargin=1.4em,itemsep=0.3em]
\item \textbf{Rouven E.\ Haschka}, \emph{Copula-based endogeneity corrections
      in R}, \\
      \url{https://github.com/HashtagHaschka/Copula-based-endogeneity-corrections} \\
      ORCID \url{https://orcid.org/0000-0002-2916-9745}. This is the reference
      implementation. Every estimator, every CDF estimator, the bootstrap, the
      diagnostics and the validity report are ported from it.
\item \textbf{Ashwin Malshe}, \emph{endogCopula: Endogeneity Corrections via
      Gaussian Copulas}, \\
      \url{https://github.com/ashgreat/endogCopula}. The packaged R version of
      the same estimators, which shaped how the port is organised into modules
      and what the documented interface looks like.
\end{itemize}

If this code contributes to published work, please cite both of those, and the
paper behind whichever estimator was used. The list is in
Section~\ref{sec:refs}.

\section{The model}

Let $y$ be the outcome, $P$ a vector of endogenous regressors and $W$ a vector
of exogenous ones. The structural model is
\begin{equation}
y_i \;=\; \mu + P_i'\alpha + W_i'\beta + \xi_i ,
\qquad \operatorname{corr}(P_i, \xi_i) \neq 0 ,
\label{eq:structural}
\end{equation}
so OLS is biased. With an instrument one would run two-stage least squares.
The copula approach replaces the instrument with a distributional assumption:
the dependence between $P$ and $\xi$ is a Gaussian copula. Writing
$P^*_k = \Phi^{-1}(F_k(P_k))$ for the normal score of the $k$th endogenous
regressor and $\xi^* = \xi/\sigma$, the assumption is that $(P^*, \xi^*)$ is
jointly normal. Then
\begin{equation}
E[\xi_i \mid P_i] \;=\; \sigma \sum_k \rho_k P^*_{ik},
\label{eq:cf}
\end{equation}
and adding $\hat P^*$ to \eqref{eq:structural} as a control function removes
the bias. The estimator is ordinary least squares on the augmented regression
\begin{equation}
y_i \;=\; \mu + P_i'\alpha + W_i'\beta + C_i'\gamma + u_i .
\label{eq:augmented}
\end{equation}

What identifies the model is \emph{non-normality} of $P$. If $P$ is normal
then $\Phi^{-1}(F(P))$ is an affine function of $P$ itself, $C$ is collinear
with the regressors, and \eqref{eq:augmented} has no solution. This is not a
technicality to be checked afterwards: it is the entire source of
identification, and it is why the \stcmd{validity} report leads with the
shape of $P$.

In a finite sample the failure is quieter than that, and worth recognising.
$\hat F$ is not the true normal CDF, so a near-normal $P$ leaves a design that
is nearly rather than exactly singular: the model estimates, and the estimate
is badly wrong. On $n = 400$ draws with a standard normal $P$ and a true
$\alpha$ of $2.0$, the correction returns $2.54$ with a standard error of
$0.51$ against an OLS standard error of $0.03$ --- an ICON above 15. Nothing
errors out. The diagnostics are what catch it, which is why they are printed
rather than merely stored.

The five estimators differ in how $C$ is built.

\subsection{The estimators}

\textbf{PG, Park and Gupta (2012).} The original.
$C_k = \Phi^{-1}(\hat F_k(P_k))$, added directly. Assumes
$\operatorname{corr}(W, P^*) = 0$. This assumption is what motivated
everything that followed, and it is testable: \stcmd{validity} reports it.

\textbf{2sCOPE, Yang, Qian and Xie (2025).} Transforms the exogenous
regressors too, and residualises:
\begin{equation}
C_k = P^*_k - \delta_k' W^* ,
\end{equation}
where $W^* = \Phi^{-1}(\hat F(W))$ and $\delta_k$ comes from a first-stage
regression of $P^*_k$ on $W^*$ alone --- never on the other endogenous
regressors. The correction is then orthogonal to $W$ by construction, so the
uncorrelatedness assumption is not needed. The first stage carries an
intercept.

\textbf{IMA, Haschka (2025a).} The same construction without an intercept in
the first stage, which is what the paper's derivation of $\rho$ requires: the
slope of $P^*$ on $W^*$ is then the Pearson correlation between the two. In
practice the two estimators are very close; the gap widens the further the
transformed exogenous regressors sit from mean zero.

\textbf{BMW, Breitung, Mayer and Wied (2024).} Reverses the order of the two
operations. First regress $P$ on $W$ by OLS on the \emph{raw} variables and
keep the residual $\hat e$; then take its rank transform,
$C = \Phi^{-1}(\text{rank}(\hat e)/(n+1))$. Same two ingredients, opposite
order, and the two are not reparameterisations of one another: their abstract
puts the difference plainly, that 2sCOPE ``calculate[s] residuals from a
regression of the transformed regressors, thereby assuming a particular
nonlinear dependence.'' Two consequences show up in the output. The
identification requirement now falls on $\hat e$ rather than on $P$, so
\stcmd{validity} tests $\hat e$; and their Corollary 3.2 makes the textbook
$t$ statistic valid for testing $\rho = 0$, so a Durbin--Hausman--Wu test is
reported next to the bootstrap one.

\textbf{JAMS, Liengaard et al.\ (2025).} Lets the copula structure differ
across the categories of the discrete exogenous regressors. With $d_P$
endogenous and $d_W$ continuous exogenous regressors,
\begin{equation}
C(P_i, W_i) = \bigl( C(P_i)' \;\; C(W_i)' \bigr)\,\Sigma^{-1}
              \begin{bmatrix} I_{d_P} \\ 0 \end{bmatrix},
\end{equation}
with $\Sigma$ the variance--covariance matrix of $(C(P), C(W))$ --- the
covariance matrix, not the correlation matrix; the two give different linear
combinations and therefore different estimates. Conditioning on a set of
categorical controls estimates $\Sigma$ and the marginal CDFs inside each joint
category, so every copula column is zero outside its own cell.

\subsection{Estimating the marginal CDF}

The literature disagrees on how $\hat F$ should be estimated, so all of the
proposals are implemented and the choice is free (except in BMW).

\begin{center}
\small
\begin{tabular}{llp{6.4cm}}
\toprule
Option & Source & What it is \\
\midrule
\stcmd{kde.silverman} & Park \& Gupta (2012)  & integrated Epanechnikov kernel density, Silverman bandwidth \\
\stcmd{kde.cv}        & Li, Li \& Racine (2017) & the same, cross-validated bandwidth \\
\stcmd{kde.plugin}    & Polansky \& Baker (2000) & Gaussian kernel CDF, plug-in bandwidth \\
\stcmd{ecdf.fixed}    & Becker et al.\ (2022) & empirical CDF with replaced boundary \\
\stcmd{ecdf.adj}      & Liengaard et al.\ (2025) & adjusted ECDF \\
\stcmd{rank.n}        & Qian, Koschmann \& Xie (2025) & rank$/n$ with a top correction \\
\stcmd{rank.n1}       & Breitung et al.\ (2024) & rank$/(n{+}1)$ \\
\bottomrule
\end{tabular}
\end{center}

Defaults follow each estimator's own paper: \stcmd{kde.silverman} for PG,
\stcmd{rank.n} for 2sCOPE and IMA, \stcmd{ecdf.adj} for JAMS, \stcmd{rank.n1}
for BMW. BMW is the one place where the choice is not free, because its
Proposition 3.1 is derived for rank$/(n{+}1)$.

\stcmd{kde.cv} is available in Python but not in Stata: it needs an $O(n^2)$
cross-validation inside every bootstrap replicate.

\subsection{Inference}

$C$ is a generated regressor, so the textbook OLS standard errors from
\eqref{eq:augmented} are wrong for the structural coefficients. Inference is
by a plain pairs bootstrap: draw row indices with replacement, recompute the
copula terms on the resample, refit. This is the procedure used in every one
of the underlying papers. Resamples that lose a factor level, or that leave a
JAMS category too thin to estimate its own $\Sigma$, are rejected and redrawn;
if that happens often the command says so, because the standard errors are
then conditional on the draws that survived.

Each fit also carries the uncorrected OLS estimate from the \emph{same}
resamples. The ratio of the two standard errors is ICON, the standard error
inflation of Qian, Koschmann and Xie (2025). Values above 6 flag weak
identification or a misspecified dependence model.

\section{The Stata command}

\subsection{Syntax}

\begin{stata}
copulaendog depvar endogvars [if] [in] [, exog(varlist) method(name)
    cdf(name) ties(name) nboots(#) seed(#) conditional(varlist)
    discrete(varlist) fsexclude(varlist) generate(stub) noconstant
    level(#) validity ]
\end{stata}

Position decides. Variables listed after \emph{depvar} are endogenous and each
gets its own copula term; everything in \stcmd{exog()} is exogenous and gets
none. Endogenous regressors must be plain numeric variables, because a factor
variable expands into several columns and has no single copula term; put
categorical controls in \stcmd{exog()}, where factor-variable notation is
allowed.

\begin{center}
\small
\begin{tabular}{lll}
\toprule
Option & Default & Meaning \\
\midrule
\stcmd{exog()}        & --- & exogenous regressors, factor variables allowed \\
\stcmd{method()}      & \stcmd{pg} & \stcmd{pg}, \stcmd{2scope}, \stcmd{ima}, \stcmd{bmw}, \stcmd{jams} \\
\stcmd{cdf()}         & per estimator & marginal CDF estimator \\
\stcmd{ties()}        & \stcmd{max} & \stcmd{max} for the counting function, \stcmd{average} for midranks \\
\stcmd{nboots()}      & 199 & bootstrap replicates \\
\stcmd{seed()}        & --- & seed; without it the draws are not reproducible \\
\stcmd{conditional()} & --- & JAMS: what the copula structure may vary over \\
\stcmd{discrete()}    & --- & JAMS: exogenous regressors to keep out of $C$ \\
\stcmd{fsexclude()}   & --- & exogenous regressors to hold out of the first stage \\
\stcmd{generate()}    & --- & keep the copula terms as variables \\
\stcmd{validity}      & --- & report the identification checks \\
\bottomrule
\end{tabular}
\end{center}

\subsection{A worked example}

The data are simulated so that the assumptions can be checked against the
truth. The endogenous regressor is $P = \exp(0.8 v + 0.6 w)$ with
$\operatorname{corr}(v, \xi) = 0.6$, and
\[
y = 1 + 2P + 1.5x - 0.5w + \xi .
\]
Because $P$ is a strictly monotone function of the normal variate
$0.8v + 0.6w$, the Gaussian copula between $P$ and $\xi$ holds exactly, while
$P$ itself is lognormal and therefore strongly non-normal --- exactly the
setting the estimators are built for. Note the $0.6w$ term: it makes
$\operatorname{corr}(W, P^*) \neq 0$, so the Park and Gupta assumption fails by
construction and the later estimators should do better.

\begin{stata}
. use simdata, clear
. regress y P x w
. copulaendog y P, exog(x w) method(pg)     nboots(499) seed(1) validity
. copulaendog y P, exog(x w) method(2scope) nboots(499) seed(1)
. copulaendog y P, exog(x w) method(bmw)    nboots(499) seed(1)
\end{stata}

On $n = 2000$ draws from that design the estimates of $\alpha$, whose true
value is $2.000$, are as follows. (These are from the Python port, which has
been checked against the R reference to $10^{-12}$; the file is
\texttt{validation/vignette\_data.csv}, so they can be reproduced in any of the
three languages.)

\begin{center}
\small
\begin{tabular}{lrrrr}
\toprule
& $\hat\alpha$ & boot.\ SE & $\hat\rho$ & max ICON \\
\midrule
OLS, uncorrected & 2.2066 & 0.0106 & --- & --- \\
PG               & 2.0039 & 0.0320 & 0.669 & 0.91 \\
2sCOPE           & 2.0136 & 0.0258 & 0.470 & 1.12 \\
IMA              & 2.0136 & 0.0258 & 0.470 & 1.12 \\
BMW              & 2.0109 & 0.0220 & 0.512 & 1.54 \\
JAMS             & 2.0125 & 0.0254 & 0.470 & 1.13 \\
\bottomrule
\end{tabular}
\end{center}

Uncorrected OLS is off by $0.21$, about twenty of its own standard errors. All
five corrections land within $0.014$ of the truth, and every ICON is close to
one, so the correction costs almost nothing in precision here. The price of
the correction is visible all the same: the standard error on $\hat\alpha$
triples under PG.

The \stcmd{validity} report on the same fit rejects the uncorrelatedness
assumption --- correctly, since we built the data that way, with
$R^2 = 0.349$, $F = 535.1$ and $p < 10^{-100}$ for the joint test --- and
points the user at the estimators that do not need it. It also confirms that
$P$ is non-normal enough to identify the model on both the Yang and the Becker
criteria (skewness $10.0$, AD $182.9$, KS $p < 10^{-90}$), which is what makes
this design work at all.

Two endogenous regressors, a factor control, and JAMS with the copula
structure varying over the categories of \stcmd{store}:

\begin{stata}
. copulaendog y price adspend, exog(feature) method(2scope) seed(1)
. copulaendog y price, exog(feature i.store) method(bmw) seed(1)
. copulaendog y price, exog(feature i.store) method(jams) ///
      conditional(store) seed(1)
\end{stata}

Prediction excludes the copula terms, because they are endogeneity controls
and not part of the causal model:

\begin{stata}
. predict yhat
. predict xi, residuals
\end{stata}

\subsection{Reading the validity report}

The report keeps three sources apart rather than merging them into one verdict,
because they do not fully agree.

\begin{enumerate}[leftmargin=1.6em,itemsep=0.3em]
\item \textbf{Non-normality} of $P$ --- or of the first-stage residuals under
      BMW. Two criteria are shown. ``Yang ok'' is the Kolmogorov--Smirnov
      criterion of Yang, Qian and Xie (2025), $p < .05$, which they prefer
      precisely because it is conservative. ``Becker ok'' is the C5.0 decision
      tree of Becker, Proksch and Ringle (2022, Fig.~8), which combines
      skewness with the Anderson--Darling and Cram\'er--von Mises statistics
      and depends on the sample size. Becker et al.\ report a correlation of
      $.66$ between the AD statistic and the copula term's $t$ statistic
      against $-.03$ for KS, which is why both are shown.
\item \textbf{The uncorrelatedness assumption}, printed only for PG. Both the
      Holm-adjusted single correlations and a joint test are reported; the
      assumption is about the whole set, so the joint test is the one that
      matters.
\item \textbf{The structural error.} Becker et al.\ find that a non-normal
      error breaks consistency in models with an intercept; Yang et al.\ and
      Qian et al.\ permit one under the decomposition $\xi = U + V$. The
      residuals show $\xi$ rather than $U$, so their shape settles neither
      question, and the report says so.
\item \textbf{ICON}, the standard error inflation relative to uncorrected OLS.
\end{enumerate}

\section{The Python package}

The same five estimators, the same seven CDF estimators, the same bootstrap:

\begin{stata}
import pandas as pd
from copulaendog import copreg, sim_endog

df  = sim_endog(n=2000, seed=123)
fit = copreg("y ~ P | x + w", df, method="2scope", nboots=499, seed=1)

print(fit.summary())
print(fit.validity())

fit.params, fit.bse, fit.vcov      # coefficients and bootstrap covariance
fit.conf_int(0.95, kind="percentile")
fit.rho, fit.icon                  # endogeneity measure, SE inflation
fit.predict(newdata)               # structural model only
\end{stata}

The formula follows the R interface: terms before the \verb=|= are endogenous,
terms after it are exogenous. Everything \texttt{patsy} accepts works ---
\verb=np.log(price)=, \verb=price:feat=, \verb=I(x**2)=, \verb=C(store)=, and
\verb=- 1= to drop the intercept.

Individual estimators can be called directly --- \texttt{CopRegPG},
\texttt{CopReg2sCOPE}, \texttt{CopRegIMA}, \texttt{CopRegBMW},
\texttt{CopRegJAMS} --- and the copula transformation itself is exposed as
\texttt{copula\_transform(x, cdf, ties)} for use in models this package does
not cover.

\section{What is not ported}

Three estimators from the R implementation are deliberately out of scope here,
and the R code should be used for them:

\begin{itemize}[leftmargin=1.4em,itemsep=0.2em]
\item \textbf{2sCOPE-np} (Hu, Qian and Xie 2025), which replaces the first
      stage with a nonparametric conditional CDF estimated by the kernel
      method of Li and Racine (2008). It is the only estimator in the family
      that tolerates discrete endogenous regressors.
\item \textbf{PANEL} (Haschka 2022), a linear panel model with fixed effects
      estimated by maximum likelihood after a forward orthogonal deviations
      transformation.
\item \textbf{BAYES} (Haschka 2025b), in which the marginal CDFs are unknown
      parameters drawn alongside the coefficients rather than plugged in.
\end{itemize}

\section{Verification}

The Python port is checked against the R reference implementation on a common
dataset: across ten model specifications --- all five estimators, four CDF
estimators, an interaction, and JAMS with and without conditioning --- all 75
coefficients and $\rho$ values agree to within $10^{-12}$, and the six
closed-form CDF estimators agree to machine precision. The comparison script
is in \texttt{validation/} and can be rerun.

The Stata port implements the identical algorithms in Mata and ships with a
self-test do-file, \texttt{validation/stata\_selftest.do}, which runs every
estimator, compares the copula transformations against the same R reference,
and prints the coefficients for comparison. Run it once in your own Stata
installation before relying on the results.

\section{References}
\label{sec:refs}

\small
\begin{list}{}{\leftmargin=1.5em \itemindent=-1.5em \itemsep=0.25em}

\item Becker, J.-M., D. Proksch, and C. M. Ringle. 2022. Revisiting Gaussian
copulas to handle endogenous regressors. \emph{Journal of the Academy of
Marketing Science} 50: 46--66.

\item Breitung, J., A. Mayer, and D. Wied. 2024. Asymptotic properties of
endogeneity corrections using nonlinear transformations. \emph{The
Econometrics Journal} 27(3): 362--383.

\item Haschka, R. E. 2022. Handling endogenous regressors using copulas: A
generalization to linear panel models with fixed effects and correlated
regressors. \emph{Journal of Marketing Research} 59(4): 861--880.

\item Haschka, R. E. 2025a. Robustness of copula-correction models in causal
analysis: Exploiting between-regressor correlation. \emph{IMA Journal of
Management Mathematics} 36(1): 161--180.

\item Haschka, R. E. 2025b. Bayesian inference for joint estimation models
using copulas to handle endogenous regressors. \emph{Oxford Bulletin of
Economics and Statistics}.

\item Hu, X., Y. Qian, and H. Xie. 2025. Correcting endogeneity via
nonparametric copula control functions. NBER Working Paper 33607.

\item Li, Q., and J. S. Racine. 2008. Nonparametric estimation of conditional
CDF and quantile functions with mixed categorical and continuous data.
\emph{Journal of Business \& Economic Statistics} 26(4): 423--434.

\item Li, Q., J. Li, and J. S. Racine. 2017. Cross-validated mixed-datatype
bandwidth selection for nonparametric cumulative distribution/survivor
functions. \emph{Econometric Reviews} 36: 970--987.

\item Liengaard, B. D., J.-M. Becker, M. Bennedsen, P. Heiler, L. N. Taylor,
and C. M. Ringle. 2025. Dealing with regression models' endogeneity by means
of an adjusted estimator for the Gaussian copula approach. \emph{Journal of
the Academy of Marketing Science} 53: 279--299.

\item Park, S., and S. Gupta. 2012. Handling endogenous regressors by joint
estimation using copulas. \emph{Marketing Science} 31(4): 567--586.

\item Polansky, A. M., and E. R. Baker. 2000. Multistage plug-in bandwidth
selection for kernel distribution function estimates. \emph{Journal of
Statistical Computation and Simulation} 65: 63--80.

\item Qian, Y., A. Koschmann, and H. Xie. 2025. A practical guide to
endogeneity correction using copulas. \emph{Journal of Marketing}.

\item Yang, F., Y. Qian, and H. Xie. 2025. Addressing endogeneity using a
two-stage copula generated regressor approach. \emph{Journal of Marketing
Research} 62(4): 601--623.

\end{list}

\end{document}
